\subsection{正元素}

设 E 为复 Hilbert 空间，u 为 E 的自同态。回忆（EVT，V，第 45 页，定义 6）：若 \(\langle x|u(x)\rangle\geqslant0\) 对所有 \(x\in E\) 成立，则称 u 是正的。这时自同态 u 是 Hermite 的（同前）。此外，若 F 是复 Hilbert 空间且 \(v\in\mathcal{L}(\mathrm{F};\mathrm{E})\) ，则 F 的自同态 \(v^{*}uv\) 是正的（EVT，V，第 45 页，命题 12）。

命题 8.　设 E 为复 Hilbert 空间，u 为 E 的自同态。下列条件等价：

\begin{enumerate}
\def\labelenumi{(\roman{enumi})}
\item
  自同态 u 是正的；
\item
  u 的数值域包含于 \(\mathbf{R}_{+}\) ；
\item
  自同态 \(u\) 是星代数 \(\mathcal{L}(\mathrm{E})\) 的正元素；
\item
  存在 \(\mathcal{L}(\mathrm{E})\) 的 Hermite 元素 v，使得 \(u = v^2\) ；
\item
  存在从 E 到复 Hilbert 空间 F 中的连续线性映射 v，使得 \(u = v^{*}v\) 。
\end{enumerate}

根据 EVT，V，第 45 页，定义 6，u 是正的当且仅当它是 Hermite 的且 \(\langle x|u(x)\rangle\geqslant0\) 对所有 \(x\in E\) 成立。(i) 因此，蕴含关系 (i) \(\Longrightarrow\) (ii) 由数值域的定义得出。

\begin{enumerate}
\def\labelenumi{(\roman{enumi})}
\setcounter{enumi}{1}
\item
  \(\Longrightarrow\) (iii)：假设蕴涵 u 是 Hermite 的（EVT，V，第 45 页，及第 2 页的注记），且它的谱包含于 \(R_{+}\) （命题 6）；因此，u 是星代数 \(\mathcal{L}(\mathrm{E})\) 的正元素。
\item
  \(\Longrightarrow\) (iv)：这是 I，第 118 页，命题 16 的一个特殊情形。
\item
  \(\Longrightarrow\) (v) 是显然的。
\item
  \(\Longrightarrow\) (i)：设 F 为复 Hilbert 空间，\(v \in \mathcal{L}(\mathrm{E};\mathrm{F})\) 满足 \(u = v^{*}v\) 。设 \(x \in E\) 。我们有 \(\langle x|u(x)\rangle = \langle x|(v^{*}v)(x)\rangle = \|v(x)\|^{2}\) ，这就证明 u 是正的。
\end{enumerate}

回忆（EVT，V，第 45 页，注记 1）：对 \(\mathcal{L}(\mathrm{E})\) 的每个 Hermite 元素 u，我们置

\[
m(u) = \inf_{\substack{x\in \mathrm{E}\\ \| x\| = 1}}\langle x|u(x)\rangle = \inf \iota (u) = \inf_{x\in \operatorname{E}\text{-}\{0\}}\frac{\langle x|u(x)\rangle}{\|x\|^ {2}},
\]

\[
\operatorname{M}(u) = \sup_{\substack{x\in \operatorname{E}\\ \| x\| = 1}}\langle x|u(x)\rangle = \sup \iota (u) = \sup_{x\in \operatorname{E}\text{-}\{0\}}\frac{\langle x|u(x)\rangle}{\|x\|^ {2}}.
\]

\[
\text{若 } \mathrm{E} = \{0\}\text{，则 } \mathrm{M}(u) = -\infty,\ m(u) = +\infty\text{，且 } \iota(u) = \varnothing.
\]

假设 E 非零；这时我们有 \(m(u) \leqslant \mathrm{M}(u)\) ，且 u 的数值域是以 \(m(u)\) 与 \(\mathrm{M}(u)\) 为端点的区间。由命题 6，\(\mathrm{Sp}(u)\) 包含于区间 \([m(u), \mathrm{M}(u)]\) 。更确切地说：

命题 9.　设 E 为复 Hilbert 空间，u 为 \(\mathcal{L}(\mathrm{E})\) 的 Hermite 元素。

\begin{enumerate}
\def\labelenumi{\alph{enumi})}
\item
  我们有 \(m(u) = \inf \operatorname{Sp}(u)\) 与 \(\operatorname{M}(u) = \sup \operatorname{Sp}(u)\) ；
\item
  若 E 非零，我们有 \(\|u\|=\sup(|m(u)|,|\mathrm{M}(u)|)\) 。
\end{enumerate}

设 \(\lambda \in \mathbf{R}\) 。\(\lambda\) 是 \(u\) 的谱的下界的充分必要条件是 \(u - \lambda \geqslant 0\) 。这（命题 8，(ii)）等价于条件 \(\langle x|u(x)\rangle \geqslant \lambda \|x\|^2\) 对所有 \(x \in E\) 成立，即等价于 \(m(u) \geqslant \lambda\) 。这就证明 \(m(u)\) 是 \(\mathrm{Sp}(u)\) 的下确界。类似地，可以验证 \(\mathrm{M}(u)\) 是 \(\mathrm{Sp}(u)\) 的上确界。

由于 u 是正规的，我们有 \(\varrho(u)=\|u\|\) （I，第 108 页，推论 1）。由于 \(E\neq\{0\}\) ，u 的谱非空（I，第 26 页，推论 1），且 \(\varrho(u)\) 是包含 \(\mathrm{Sp}(u)\) 的以 0 为心的最小圆盘的半径（I，第 24 页，定理 1），因此 b) 由 a) 得出。

